\scanpage{103}
\mspage{2}
% archive foliation: 102
\[
\underline\omega+\underline O_Xs_0+\underline\omega^{-1}\nu_{s_0}
\]
\[
s=s_0+a_1=(\exp\Theta_{-a_1})s_0=s_0+[-a_1,s_0].
\]
\[
\left[\begin{aligned}
\nu_s&=(\exp\Theta_{-a_1}\otimes\mathrm{id}_{\underline\omega})\nu_{s_0}
=1+\underbrace{[-a_1,\nu_{s_0}]}_{-2a_1\unsure{\nu_{s_0}}}
+\frac12\underbrace{[-a_1,\underbrace{[-a_1,\nu_{s_0}]}_{-2a_1}]}_{-\unsure{2}a_1^2}\\
\nu_s&=\deleted{\unsure{a}}\underbrace{-a_1^2}_{2}\underbrace{-2a_1}_{1}+\underbrace{1}_{0}
\end{aligned}\right]
\]
Il faut calculer, pour avoir l'équation d'horizontalité sur $s$ (ou sur $a_1$) relative à $c$, la \deleted{\unsure{coefficient}}\inserted{composante} dans $\underline\omega^2$ (première comp.) de
\[
(\exp(\Theta_{a_1})\otimes\mathrm{id}_{\underline\omega})(c(s))
\]
et l'égaler à 0.

\deleted{On \unsure{supp. $c$ provenant d'une conn. $c$ associée} à un splitting, i.e. pour lequel ce splitting est horizontal}
\deleted{$c(s)=c(s_0)+c(a_1)$}
\[
c(s_0)=\deleted{\unsure{\bar\varphi_1+\nu_{s_0}}}\bar\varphi_2+\nu_{s_0}
\]
\[
(\exp\Theta_{a_1}\otimes\mathrm{id}_{\underline\omega})
\underbrace{(c(s))}_{\substack{c(s_0)+c(a_1)\\=\bar\omega_2+c(a_1)+\nu_{s_0}}}
\quad\deleted{\unsure{+}}(\exp\Theta_{a_1}\otimes\mathrm{id}_{\underline\omega})(\unsure{a_1})+{}
\]
\unsure{est sa comp.} dans $\underline\omega^2$, et si
\[
c(a_1)=\deleted{\unsure{\omega}}\underbrace{\overbrace{D_1(a_1)}^{2}}_{\substack{\text{op. diff. d'ordre}\\1\quad\underline\omega\longrightarrow\underline\omega^2}}
-\overbrace{2a_1}^{1}+\deleted{\unsure{\nu_{s_0}}}\ (\deleted{\unsure{\exp\Theta\omega}}\ldots
\]
\margin{pour un splitting convenable \inserted{et une base $\delta$ de $\omega$}, \deleted{\unsure{si on a}} on a $D_1(\delta)=0$ i.e. $D_1(\lambda\delta)=\delta\,d\lambda$}
On a donc
\[
\begin{aligned}
(\exp(\Theta_{a_1})\otimes\mathrm{id}_{\underline\omega})(c(s))
&=\inserted{\underline\omega_2+{}}\overbrace{D_1(a_1)}^{2}-\overbrace{2a_1}^{1}
-\underbrace{\overbrace{2[a_1,a_1]}^{2}}_{+2a_1a_1}\\
&\qquad\deleted{+2\unsure{a}}\underbrace{-a_1^2}_{2}+2a_1+1
\end{aligned}
\]
L'équation est donc
\[
\boxed{D_1(a_1)+a_1^2\unsure{-|+}\underline\omega_2=0}
\]
\margin{\unsure{facteur 2 qui manque avant ce calcul ?}}
\margin{équation d'horizontalité de la section de $P$ \unsure{disjointe} de $\sigma$, définie par la 1-forme $a_1$}
Si on choisit convenablement une base de $\omega$ (p. ex. $\delta=dt$) on peut supposer $D_1(\delta)=0$ et l'équation en termes de $a_1=\lambda\delta$ est
\[
(d\lambda)\delta+\lambda^2\delta^2=0
\]
ou encore, si $\delta=dt$ i.e. $a_1=\lambda\,dt$
\deleted{$\unsure{\frac{d\lambda}{dt}+\lambda}$}
\[
\boxed{\frac{d\lambda}{dt}+\lambda^2=A}
\]
\margin{TSVP}
